\subsection{Graded regular algebras}

In this issue, we assume that \(H_{0}\) is a field and that H is an \(H_{0}\)-algebra of finite type.

One sets \(H_{+} = H_{\geqslant 1}\) ; it is a maximal ideal of H. The ring \(S^{-1}H\) is identified with the local ring \(H_{H_{+}}\) of H at the ideal \(H_{+}\) ; its maximal ideal is \((\mathrm{H}_{+})_{\mathrm{H}_{+}} = \mathrm{S}^{-1}\mathrm{H}_{+}\) , its residue field is identified with \(H_{0}\) .

PROPOSITION 8. --- a) We have \(\dim(\mathbf{H}) \leqslant [\mathbf{H}_{+}/\mathbf{H}_{+}^{2}:\mathbf{H}_{0}]\) .

\begin{enumerate}
\def\labelenumi{\alph{enumi})}
\setcounter{enumi}{1}
\tightlist
\item
  The following conditions are equivalent:
\end{enumerate}

\begin{enumerate}
\def\labelenumi{(\roman{enumi})}
\item
  \(\dim (\mathrm{H}) = [\mathrm{H}_{+} / \mathrm{H}_{+}^{2}:\mathrm{H}_{0}]\) ;
\item
  the Noetherian local ring \(\mathbf{S}^{-1}\mathbf{H}\) is regular ;
\item
  H is generated as \(H_{0}\) -algebra by a family of homogeneous elements of degrees \textgreater{} 0, algebraically independent over \(H_{0}\) .
\end{enumerate}

\begin{enumerate}
\def\labelenumi{\alph{enumi})}
\setcounter{enumi}{2}
\tightlist
\item
  Suppose the conditions of b) are satisfied, and let \(a_{1}, \ldots, a_{n} \in \mathbf{H}\) be homogeneous elements of degrees \textgreater{} 0. The following conditions are equivalent:
\end{enumerate}

\begin{enumerate}
\def\labelenumi{(\roman{enumi})}
\item
  the classes of \(a_{i}\) form a basis of the \(\mathbf{H}_{0}\) -vector space \(\mathbf{H}_{+}/\mathbf{H}_{+}^{2}\) ;
\item
  the \(a_{i}/1\) form a coordinate system of the regular Noetherian local ring \(\mathbf{S}^{-1}\mathbf{H}\) ;
\item
  the \(a_{i}\) are algebraically independent over \(\mathbf{H}_{0}\) and generate \(\mathbf{H}\) as \(\mathbf{H}_{0}\) -algebra.
\end{enumerate}

We have \(\dim(\mathbf{H})=\dim(\mathbf{S}^{-1}\mathbf{H})\) (no. 2, th. 1), and \([\mathbf{H}_{+}/\mathbf{H}_{+}^{2}:\mathbf{H}_{0}]=[(\mathbf{S}^{-1}\mathbf{H}_{+})/(\mathbf{S}^{-1}\mathbf{H}_{+})^{2}:\mathbf{H}_{0}]\) (II, § 3, no. 3, prop. 9); according to § 5, no. 1, this implies \(a\) ) and the equivalences (i) \(\Leftrightarrow\) (ii) in \(b\) ) and \(c\) ). The two implications (iii) \(\Rightarrow\) (i) in \(b\) ) and \(c\) ) are trivial. Let us prove the implications (i) \(\Rightarrow\) (iii). Suppose therefore that we have \(\dim(\mathbf{H})=[\mathbf{H}_{+}/\mathbf{H}_{+}^{2}:\mathbf{H}_{0}]\) and let \(a_{1},\ldots,a_{n}\) be homogeneous elements of H, of degrees \(>0\) , whose classes form a basis of \(\mathbf{H}_{+}/\mathbf{H}_{+}^{2}\) ; consider the homomorphism of graded algebras \(\varphi:\mathrm{H}_{0}[\mathrm{X}_{1},\ldots,\mathrm{X}_{n}]\to\mathrm{H}\) which sends \(\mathrm{X}_{i}\) onto \(a_{i}\) . The ideal \(\mathrm{H}_{+}\) of H is generated by the \(a_{i}\) (A, II, p.~171, cor. 2 and remark); consequently the homomorphism \(\varphi\) is surjective (III, § 1, no. 2, prop. 1). But we have

\[
\dim \left(\mathrm{H} _ {0} \left[ \mathrm{X} _ {1}, \dots , \mathrm{X} _ {n} \right]\right) = n = \dim (\mathrm{H}) = \dim \left(\mathrm{H} _ {0} \left[ \mathrm{X} _ {1}, \dots , \mathrm{X} _ {n} \right] / (\operatorname{Ker} \varphi)\right)
\]

(§ 2, no. 4, cor. 1 of th. 3), hence Ker φ = 0 since H₀{[}X₁, \ldots, Xₙ{]} is an integral domain. This implies assertions (iii).

Under the hypotheses of \(b\) ), one says that H is a regular \(\mathrm{H}_{0}\) -graded algebra, or a \(\mathrm{H}_{0}\) -graded polynomial algebra. Under the hypotheses of \(c\) ), we say that the \(a_{i}\) form a graded coordinate system of H.

Remarks. --- 1) With the notation of \(c\) ), let \(d_{i}\) be the degree of \(a_{i}\) ( \(1 \leqslant i \leqslant n\) ). Then the Poincaré series \(\mathbf{P}_{\mathrm{H}} = \sum_{n \in \mathbf{Z}} [\mathrm{H}_{n} : \mathrm{H}_{0}] \cdot \mathrm{T}^{n}\) is equal to \(\prod_{i} (1 - \mathrm{T}^{d_{i}})^{-1}\) (§4, no. 2, example 1). Consequently, if H is a \(\mathrm{H}_{0}\) -graded polynomial algebra, we have

\[
\mathrm{P} _ {\mathrm{H}} = \prod_ {n} (1 - \mathrm{T} ^ {n}) ^ {- \delta (n)}, \quad \text { with } \quad \delta (n) = \left[ (\mathrm{H} _ {+} / \mathrm{H} _ {+} ^ {2}) _ {n}: \mathrm{H} _ {0} \right].
\]

\begin{enumerate}
\def\labelenumi{\arabic{enumi})}
\setcounter{enumi}{1}
\tightlist
\item
  Conversely, the fact that there exist integers \(d_{i}>0\) such that \(\mathbf{P}_{\mathbf{H}}=\prod_{i}(1-\mathbf{T}^{d_{i}})^{-1}\) does not imply that H is a graded polynomial algebra; for example, if H is generated by an element X of degree 1 and an element Y of degree 2, subject only to the relation \(\mathbf{X}^{2}=0\) , one has \(\mathbf{P}_{\mathbf{H}}=(1-\mathbf{T})^{-1}\) .
\end{enumerate}

